\section{Normalized grid as the reference grid}
In this section it is described which grid has been implemented in the new
version of the code. The form that the transport equations acquire when
written in terms of this grid is also presented and discussed. 
%
\subsection{Motivations for the change of grid}
In Astra--7.0 it has been decided to adopt another grid, which has
advantages and disadvantages over the grid used in previous versions.
The motivations for this choice stem from three considerations:

1) With the previous, $\rho$--grid, the number of grid points in the plasma region
varies with the plasma linear dimension. This means that simulations which
involve expanding plasmas with much different volume at the beginning and at the
end of the expansion require a large number of grid points to maintain a sufficient
resolution at every time step;

2) The boundary point does not fall on a pre--existing grid point,
meaning that the last two grid points acquire a differential $\delta\rho$ which varies with time;

3) Coupling with equilibrium codes is extremely difficult and prone to numerical instabilities.

The main advantage of the $\rho$--grid is that grid--convective terms appear 
only for cases with 
time--dependent vacuum field $\dot{B}_0$. In case of static vacuum field, 
changes in plasma boundary do not affect internal grid points, which 
simplifies the equations. However, in the general case, the equations have the
same form as they would appear on the normalized grid, where in the latter case $\dot{\Phi}_{\rm b}$ appears as the convective drive.
%
\subsection{Description of the chosen grid}
From this version onwards, the fundamental grid on which quantities are computed is defined as:
\begin{eqnarray}
x = \sqrt{\frac{\Phi}{\Phi_{\rm b}}}
\label{normgridfeeq}
\end{eqnarray}
In terms of indexing, the main transport grid is numerically defined as ($j  \in [1..{\rm NA1}]$):
\begin{eqnarray}
x_j = \frac{(j-0.5)}{({\rm NA1}-0.5)}
\label{normgridindexeq}
\end{eqnarray}
where ${\rm  NA1}$ is the number of grid points in the plasma. The grid is thus equispaced.
As such we have that $x_{\rm NA1}=1$ is always a grid point. 
The shifted grid on which fluxes and gradients are defined is given as:
\begin{eqnarray}
\tilde{x}_j = \frac{(j-1.)}{({\rm NA1}-0.5)}
\label{snormgridindexeq}
\end{eqnarray}
This choice of indexing is useful if simulations are performed at given boundary
values of the main quantity. For flux--driven simulations, it would be more
appropriate and accurate to have $\tilde{x}_{\rm  NA1}=1.$. 

To allow the user to choose between the two possibilities, i.e. having the 
actual plasma boundary on the main or on the shifted grid, an option has been 
introduced. The variable {\bf FLXDR} can now be set to 0 or to 1 to this purpose, as seen in table 1.1 below. 
\\
{\bf Table 1.1} Choice for grid indexing label {\bf FLXDR}\\[2ex]
\begin{tabular}{|l|l|l|}					\hline
\bf Name	&\bf Value & Description\\ \hline
\tt FLXDR 	& 0 & The main grid $x=1$ coincides with the plasma boundary\\
\tt  	& 1 &  The shifted grid $\tilde{x}=1$ coincides with the plasma boundary\\ 
\hline
\end{tabular}
\bigskip\\
The variable {\bf FLXDR} is a control parameter and can be added in the 'Control parameters' section of the {\it /equ/log/[model\_file\_name]}.
If not provided explicitly, the default value is {\bf FLXDR}=0.

{\bf IMPORTANT NOTE}: A major difference with respect to the previous Astra--6.2.1 version is that now
the number of grid points in the plasma, {\bf NA1}, has to be given in input,
instead of {\rm NB1}, the latter being then a time--dependent quantity. In 
practice, the number of grid points in between the plasma boundary and the wall 
is a varying quantity.
From the point of view of the INPUT experimental and model files to Astra, that 
is the only relevant change that the user should take care of. Every other 
detail, mapping from grid to grid etc. is taken care of inside Astra. 

The grid variables are now part of the Astra arrays set, with the names:
\\
Grid arrays\\[2ex]
\begin{tabular}{|l|l|l|}					\hline
\bf Name	& Description & Units\\ \hline
\tt XRHO 	& Main normalized grid & []\\
\tt SXHO 	& Shifted normalized grid & []\\
\tt RHO 	& Main grid in $\rho=\sqrt{\Phi/(\pi B_0)}$ & [m]\\
\tt SRHO 	& Shifted grid in $\rho=\sqrt{\Phi/(\pi B_0)}$ & [m]\\
\hline
\end{tabular}
\bigskip\\
The $\rho$ grid RHO is now defined simply as $\rho = \rho_{\rm b} x$, and SRHO as
$\tilde{\rho} = \rho_{\rm b} \tilde{x}$.
%
\subsection{Grid convection terms when working with the $x$ grid}
In the previous Astra versions, the use of $\rho$ as the fundamental grid
resulted in exact conservation of profiles for cases with vanishing transport
fluxes and sources, made exception for cases with time--dependent vacuum 
toroidal field. The latter cases were treated in terms of 'adiabatic 
compression', a posteriori (i.e. after solution of the actual transport 
equations).

When transforming equations to have time derivatives done a constant $x$ instead
of constant $\rho$, a new convective term appears, in this form:
\begin{eqnarray}
\left[\frac{\partial y}{\partial t}\right]_x - x\frac{\dot{\Phi}_{\rm b}}{2\Phi_{\rm b}} \frac{\partial y}{\partial x} + \left[{\rm transport} \right]= \left[{\rm source}\right]
\label{gridconvtermseq}
\end{eqnarray}
The second term on the LHS is the new grid convetion term, proportional to 
the quantity $\dot{\Phi}_{\rm b}$.
Because this term depends on the solution of the equilibrium problem at two 
different time slices, it is now chosen to add this convective terms {\it a 
posteriori}, as was done before, like an 'adiabatic compression/expansion'. 
The new term has also been numerically implemented in the actual matrix which 
solves for the transport equations. However then it would require iterate at 
least once between transport and equilibrium.

To allow the user to choose if to solve the convective term inside or outside 
the equations, a new variable {\bf ADCMPF} has been added to the Control 
parameters. The possible values are listed in the Table 1.2 below:
\\
{\bf Table 1.2} Choice for adiabatic compression label {\bf ADCMPF}\\[2ex]
\begin{tabular}{|l|l|l|}					\hline
\bf Name	&\bf Value & Description\\ \hline
\tt ADCMPF 	& 1 or -1 & Adiabatic compression done in the transport equation,\\ 
\tt &   &  at the moment it is done explicit (ADCMPF=1), \\
\tt &   &  improved by iterating between transport and equilibrium (NITREQ$>1$).\\
\tt &   &  Implicit scheme (ADCMPF=-1) will be implemented soon\\
\tt & 0 &  Adiabatic compression is performed outside of transport equation. \\
\tt &   &  Works only for FP! \\ 
\tt & 2 &  Adiabatic compression is neglected.\\ 
\tt &(3)&  (can then be done by equilibrium code, not recommended).\\ 
\hline
\end{tabular}
\bigskip\\
By default, {\bf ADCMPF}=1.  
We note that neglect of adiabatic compression may 
not be too dramatic for kinetic profiles in general. Exception makes the 
poloidal flux $\psi$, which suffers from severe errors if no adiabatic 
compression is computed.
%
\subsection{Choice of the numerical scheme}
In the new version of Astra, new control parameters have been added, which allow
the user to choose between several implemented numerical schemes to solve the 
transport equations.

These newly introduced variables are called {\bf INUME}{\it i}, where {\it i} can be 1,2,3, or 4. The different numbers refer to the quantities to be evolved:

1 - NE, F1, F2, ...;

2 - TE and TI;

3 - FP;

4 - UPAR;

The values that the INUME{\it i} can take are $k_1 k_2$, where $k_1$ and $k_2$ are two integer numbers. For example INUME1 = 21 means $k_1 = 2$ and $k_2=1$. 

The meaning of the $k_1,k_2$ numbers is listed and detailed in the following table 1.3:
\\
{\bf Table 1.3} Numerical scheme selection\\[2ex]
\begin{tabular}{|l|l|l|}					\hline
\bf Name	&\bf Description & Default\\ \hline
\tt $k_1$	& 1 & Scheme is fully explicit in time \\ 
\tt 	& 2 & Scheme is fully implicit in time \\ 
\tt 	& 3 & Scheme is Crank--Nicholson in time \\ 
\tt $k_2$ & 1 & Advection--diffusion is treated with centered difference in space\\
\tt  	& 2 &  The power--law scheme is employed \\ 
\tt & 3&  The exponential scheme is employed \\ 
\hline
\end{tabular}
\bigskip\\
Default value is 22 for all INUME{\it i}, meaning that by default all transport 
equations are advanced fully implicitly in time, centered difference in space
numerical schemes.
The centered difference scheme is accurate for cases with small convection and 
large diffusion. Cases with large convection are better treated with the 
exponential or the power--law scheme. In future versions, it is possible that 
more sophisticate numerical scheme with gradient limiters will be implemented, 
however in general, cases with large convection are rarely encountered.

Notice that, if {\bf INUME3}$\leq 0$, which has to be used with CU:AS; in the model file, 
then EQPF and EQFF (i.e. P' and FF') are to be prescribed by the user and are not computed in ASTRA. This
modality is useful when performing equilibrium tests at constant P', FF'.

\clearpage
